\documentclass{leveldoc}
\DocUses{foundations,deflection_plates,tune_the_lens,injection,two_stages}
\DocFingerprint{c5366cdec29942ce}

\begin{document}

\section{What the level asks}\label{sec:asks}

\leveldescription

Three positive ions leave the left edge, at $\val{aLow}$, $\val{aAxis}$ and
$\val{aHigh}$ degrees. A condenser of three apertures, a crossover gate, one
deflection plate and a fixed charge are built in. The player sets the middle
aperture's two supplies, the deflector's supply, and charges beside the fixed
charge. There are no magnets.

Each ion must pass the gate and then enter the detector moving along $+x$,
inside a cone. The gate has no direction condition
(\xeqref{two_stages}{eq:gate}). The cone is \xeqref{injection}{eq:cone}.
There is no kinetic window. A flight ends when the ion touches a charge or an
electrode, leaves the world, or enters the detector
(\xsecref{foundations}{sec:board}). Entering the detector before the gate is
a skipped gate, and the cone is not applied to that ending. An electrode's
obstacle is its box plus a contact shell of thickness $\val{contact}$. Each
charge's obstacle radius is $\val{Rcharge}$. The ion's radius is $0$.

The reference holds both halves of the middle aperture at one positive
potential and the deflector at a higher positive potential, and it adds one
positive charge. All three ions pass the gate and meet the detector's left
side inside the cone. With the supplies off and that charge omitted, none of
them does. This document uses the deflection plates' rectangle collocation.
The fixed charge and the added charge are external. There is no uniform gap.
The flight is integrated by the classical Runge--Kutta method.

\section{Setup}\label{sec:setup}

\setuptable

The speed of light is $c = \val{cLight}$. Each ion has charge $q = \val{qPart}$,
mass $m = \val{mPart}$ and radius $0$, and it is launched with kinetic energy
$T_0 = \val{T0}$ from $(\val{xLaunch}, \val{yLaunch})$. Then
$\gamma_0 = \val{gamma0}$, $\abs{\vect{p}_0} = \val{p0}$ and
$v_0 = \val{v0} = \val{vOverC}\,c$ (\xeqref{foundations}{eq:launch}).

The condenser is three apertures. The outer pairs are grounded. The reference
holds the middle pair, at $x = \val{xMid}$, at $V = \val{vMid}$. The opening
is $\val{yRimLo} \le y \le \val{yRimHi}$. The gate is
$\val{xGate0} \le x \le \val{xGate1}$, $\val{yGate0} \le y \le \val{yGate1}$.
The deflector is centred at $x = \val{xDef}$, and the reference holds it at
$V = \val{vDef}$. It occupies $\val{xDef0} \le x \le \val{xDef1}$,
$\val{yDef0} \le y \le \val{yDef1}$. The supply range is symmetric out to
$\val{vHalf}$. The fixed charge is $Q = \val{qPump}$ at
$(\val{xPump}, \val{yPump})$. The reference charge is $Q = \val{qSteer}$ at
$(\val{xSteer}, \val{ySteer})$. The largest magnitude on offer is $\val{qMag}$.

The detector is $\val{xDetLo} \le x \le \val{xDetHi}$,
$\val{yDetLo} \le y \le \val{yDetHi}$. Its cone has axis $\val{axisDeg}$
degrees and half-angle $\val{halfDeg}$ degrees. The world is
$-1 \le x \le \val{xWorld}$, $-1 \le y \le \val{yWorld}$. The reported time
step is $h = \val{H}$. The mesh uses rectangles of edge $\val{panelFine}$,
$\val{nPanels}$ of them on the upper halves of the seven electrodes.

\begin{figure}[ht]
\centering
\begin{tikzpicture}
\begin{axis}[axis equal image, width=0.92\linewidth,
  xmin=-1.5, xmax=42.5, ymin=-2, ymax=22.5,
  xlabel={$x$}, ylabel={$y$}, tick label style={font=\small}]
  \draw[black!40, dashed] (-1,-1) rectangle (\rawval{xWorld},\rawval{yWorld});
  \draw[black!25] (0,0) rectangle (40,20);
  \fill[teal!40, fill opacity=0.25]
    (\rawval{xGate0},\rawval{yGate0}) rectangle (\rawval{xGate1},\rawval{yGate1});
  \draw[teal!60!black]
    (\rawval{xGate0},\rawval{yGate0}) rectangle (\rawval{xGate1},\rawval{yGate1});
  \fill[green!50!black, fill opacity=0.18]
    (\rawval{xDetLo},\rawval{yDetLo}) rectangle (\rawval{xDetHi},\rawval{yDetHi});
  \draw[green!40!black]
    (\rawval{xDetLo},\rawval{yDetLo}) rectangle (\rawval{xDetHi},\rawval{yDetHi});
  \draw[green!50!black]
    (\rawval{coneOx},\rawval{coneOy}) -- (\rawval{coneXp},\rawval{coneYp});
  \draw[green!50!black]
    (\rawval{coneOx},\rawval{coneOy}) -- (\rawval{coneXm},\rawval{coneYm});
  \fill[black!20] (\rawval{xP0a},\rawval{yP0a}) rectangle (\rawval{xP0b},\rawval{yP0b});
  \fill[black!20] (\rawval{xP1a},\rawval{yP1a}) rectangle (\rawval{xP1b},\rawval{yP1b});
  \fill[black!20] (\rawval{xP4a},\rawval{yP4a}) rectangle (\rawval{xP4b},\rawval{yP4b});
  \fill[black!20] (\rawval{xP5a},\rawval{yP5a}) rectangle (\rawval{xP5b},\rawval{yP5b});
  \draw[orange!80!black, thick]
    (\rawval{xP2a},\rawval{yP2a}) rectangle (\rawval{xP2b},\rawval{yP2b});
  \draw[orange!80!black, thick]
    (\rawval{xP3a},\rawval{yP3a}) rectangle (\rawval{xP3b},\rawval{yP3b});
  \draw[orange!80!black, thick]
    (\rawval{xP6a},\rawval{yP6a}) rectangle (\rawval{xP6b},\rawval{yP6b});
  \fill[red!70!black] (\rawval{xPump},\rawval{yPump}) circle[radius=\rawval{Rcharge}];
  \fill[red!70!black] (\rawval{xSteer},\rawval{ySteer}) circle[radius=\rawval{Rcharge}];
  \addplot[orange!80!black] table {pathofflow.dat};
  \addplot[orange!80!black] table {pathoffaxis.dat};
  \addplot[orange!80!black] table {pathoffhigh.dat};
  \addplot[blue!60!black] table {pathlow.dat};
  \addplot[blue!60!black] table {pathaxis.dat};
  \addplot[very thick, blue!60!black] table {pathhigh.dat};
  \draw[-{Stealth[length=5pt]}, thick]
    (\rawval{xLaunch},\rawval{yLaunch}) -- (\rawval{arrowX},\rawval{arrowY});
  \fill (\rawval{xLaunch},\rawval{yLaunch}) circle[radius=0.08];
\end{axis}
\end{tikzpicture}
\caption{The column and the six flights. The outlined plates are the middle
aperture and the deflector, at the reference potentials. The gray plates are
grounded. The lower red disc is the fixed charge and the upper red disc is
the reference charge. The teal box is the gate and the green box is the
detector; the green wedge is the accepted cone. Blue paths are the reference,
and each meets the detector's left edge. Orange paths leave the supplies off
and omit the reference charge: one leaves through the right edge below the
gate, one leaves through the right edge above the launch line, and one meets
the deflector.}
\label{fig:board}
\end{figure}

\section{Principles and model}\label{sec:model}

\begin{principles}
  \item The relativistic equation of motion \xeqref{foundations}{eq:eom}, the
    launch \xeqref{foundations}{eq:launch}, and energy in a static field
    \xeqref{foundations}{eq:energy}. There is no magnetic field.
  \item Each electrode is a metal box at the potential its supply sets, or
    grounded when the supply is off (\physref{2.7}). The surface charge is
    the deflection plates' collocation
    \xeqref{deflection_plates}{eq:collocation}, with the rectangle integral
    \xeqref{deflection_plates}{eq:rectangle}. Both charges are external.
  \item The gate is \xeqref{two_stages}{eq:gate}. The detector then counts an
    entry only inside the cone \xeqref{injection}{eq:cone}.
  \item The board's rules (\xsecref{foundations}{sec:board}) and the classical
    Runge--Kutta method (\xsecref{foundations}{sec:numerics}).
  \item The radiated power \xeqref{foundations}{eq:larmor}, summed by Simpson's
    rule \xeqref{foundations}{eq:simpson}. It is not put back into the force.
\end{principles}

The reference raises three electrodes and adds one charge. The fixed charge
stays in every flight. The state obeys \xeqref{foundations}{eq:eom} with
$\vect{B} = \vect{0}$ until the first ending event. Passing the gate does not
end the flight.

\begin{assumption}{Static electrodes}\label{as:static}
Each supply is a fixed potential, with no ramp. The launch speed is
$\val{vOverC}\,c$. The game uses the same static supplies (\physref{2.7}).
\end{assumption}

\begin{assumption}{No radiation reaction}\label{as:radiation}
The radiated energy is not taken from the flight. The largest of the three
reference integrals is a fraction $\val{radFracMax}$ of $T_0$.
\end{assumption}

\begin{assumption}{No image of the ion}\label{as:image}
The ion's own surface charge is left out, as on the deflection plates and as
in the game (\physref{2.7}). The bound $q^2/d^2$, with $d$ floored at
$10^{-3}$ and taken relative to $\max(\abs{q\vect{E}}, T_0)$, reaches
$\val{refImage}$ on the reference samples, under the engine's limit
$10^{-10}$. The empty upward ray strikes the deflector, and its sampled ratio
reaches $\val{offImage}$, above that limit, because the distance falls to the
contact shell.
\end{assumption}

\begin{assumption}{Equal rectangles}\label{as:panels}
The mesh is equal rectangles of edge $\val{panelFine}$. Replacing the edge by
$\val{panelCoarse}$ moves an end state by $\val{errMesh}$. Both placements
carry electrode charge, so every ray enters that estimate.
\end{assumption}

\begin{assumption}{The fixed step}\label{as:step}
The reported step is $h = \val{H}$. Halving it estimates an end state by
$\val{errStep}$. The tolerance on a state is ten times the larger of that
estimate and the panel-size estimate, and at least $10^{-9}$. Here the
panel-size estimate is the larger one, so the tolerance is $\val{tolFlight}$.
Between samples the ion moves at most $\val{stepReach}$. The clearance of the
metal, sampled at the ends of the steps and at the ending point, reaches
$\val{refDmin}$ on the reference flights.
\end{assumption}

\section{Derivation}\label{sec:derivation}

\subsection{The condenser}\label{sec:lens}

The rectangles, the mirror and the collocation are
\xeqref{deflection_plates}{eq:rectangle} and
\xeqref{deflection_plates}{eq:collocation}. A supply replaces the potential of
the electrode it sits on. The charges' potential is subtracted on the
right-hand side. The ion is not in the linear system. The reported mesh has
$\val{nPanels}$ upper rectangles. The residual at the centroids is
$\val{solveRes}$. On the long faces, off the centroids, the potential differs
from the electrode potential by at most $\val{surfRes}$.

The middle pair is at one potential and the outer pairs are grounded, which is
the aperture lens of Tune the lens. On the launch line at the middle aperture
the transverse field is $\val{eAxis}$. Off that line, halfway to each rim, it
is $\val{eAbove}$ above and $\val{eBelow}$ below, toward the axis. The
downstream plate and the charges leave the remainder $\val{eAxis}$
(\xeqref{tune_the_lens}{eq:axis} would give zero for the lens alone). The
potential on the launch line is $\val{phiIn}$ in the first aperture,
$\val{phiMid}$ in the middle and $\val{phiOut}$ in the third, so it peaks in
the middle. \xeqref{tune_the_lens}{eq:slow} then gives $T = \val{Tmid}$ there,
the fraction $\val{slowFrac}$ of $T_0$ having gone into potential. The middle
pair carries $\val{qMid}$ and the grounded electrodes carry $\val{qOuter}$.

\subsection{The gate}\label{sec:gate}

The rule is \xeqref{two_stages}{eq:gate}. This level lists one gate and gives
it no direction condition. The flight continues through the gate. The detector
counts only after that entry. Entering the detector while the gate has not
been entered ends the flight as a skipped gate, and this document does not
then apply the cone.

\subsection{The deflector and the charges}\label{sec:drop}

There is no second plate opposite the deflector, so the uniform-gap field does
not apply. The deflector is positive. Half a cell under it the field is
$\val{eDef}$, along $-y$, so a positive ion is pushed down. The deflector
carries $\val{qDef}$. One cell above the fixed charge, that charge alone
produces $\val{eyPump}$ along $+y$, and the full field there is still
$\val{eyOver}$ along $+y$: the charge repels a positive ion upward. The
reference charge sits on the launch line, between the deflector and the
detector.

\subsection{The flight and the cone}\label{sec:flight}

The diagnostic is \xeqref{deflection_plates}{eq:flight-energy}. The state is
advanced by \xeqref{foundations}{eq:rk4} with step $h = \val{H}$. Events are
located as on the deflection plates: thirty-two Hermite samples, then sixty
bisections, obstacles before the world, and the world before the detector.
The gate is recorded when the step enters it, and the flight goes on. The
cone \xeqref{injection}{eq:cone} is applied only to a detector entry that
already passed the gate. The axis is $+x$ and the half-angle is
$\val{halfDeg}$ degrees. The step-halving estimate divides the change by $15$;
the panel-size estimate is the change itself when the edge doubles.

\section{The reference solution by hand}\label{sec:hand}

\begin{procedure}
  \item\label{st:launch} \emph{The launches} \xeqref{foundations}{eq:launch}.
    $\gamma_0 = \val{gamma0}$ and $\abs{\vect{p}_0} = \val{p0}$, so
    $v_0 = \val{v0}$. The three angles are $\val{aLow}$, $\val{aAxis}$ and
    $\val{aHigh}$ degrees.
  \item\label{st:lens} \emph{The condenser.} Hold both halves at
    $x = \val{xMid}$ at $\val{vMid}$, and leave the outer apertures grounded.
    The opening is $\val{opening}$. On the launch line,
    \xeqref{tune_the_lens}{eq:slow} gives $T = \val{Tmid}$ in the middle.
  \item\label{st:gate} \emph{The gate} \xeqref{two_stages}{eq:gate}. Each
    reference ray enters through the left face, at $x = \val{xGate0}$. The
    ordinates are $\val{LowGY}$, $\val{AxisGY}$ and $\val{HighGY}$, all
    between $\val{yGate0}$ and $\val{yGate1}$, at $t = \val{LowGT}$,
    $\val{AxisGT}$ and $\val{HighGT}$.
  \item\label{st:drop} \emph{The deflector and the charges.} Hold the
    deflector at $\val{vDef}$. Keep $Q = \val{qPump}$ at
    $(\val{xPump}, \val{yPump})$ and place $Q = \val{qSteer}$ at
    $(\val{xSteer}, \val{ySteer})$. Solve the collocation again. The centroid
    residual is $\val{solveRes}$.
  \item\label{st:march} \emph{The marches.} Advance
    \xeqref{deflection_plates}{eq:flight-energy} and
    \xeqref{foundations}{eq:eom} by \xeqref{foundations}{eq:rk4} with step
    $\val{H}$, once for each launch direction. The on-axis state every five
    time units:
    \begin{center}\small\val{states}\end{center}
  \item\label{st:end} \emph{The arrivals.} Each ion meets the left face of the
    detector. The downward ray ends at $(\val{LowX}, \val{LowY})$, at
    $t = \val{LowT}$, at an angle $\val{LowAng}$ degrees from $+x$. The
    on-axis ray ends at $(\val{AxisX}, \val{AxisY})$, at $t = \val{AxisT}$,
    at $\val{AxisAng}$ degrees. The upward ray ends at
    $(\val{HighX}, \val{HighY})$, at $t = \val{HighT}$, at $\val{HighAng}$
    degrees. The cone angles are $\val{LowBeta}$, $\val{AxisBeta}$ and
    $\val{HighBeta}$ degrees, all inside the half-angle, so each outcome is
    Arrived. The downward ray only falls, from the launch to $\val{LowY}$.
    The upward ray climbs to $\val{HighYmax}$ and then falls, and its ending
    is the lowest of the three.
  \item\label{st:check} \emph{The checks.} $W$ drifts by at most
    $\val{LowWdev}$, $\val{AxisWdev}$ and $\val{HighWdev}$. The step-halving
    estimate on the largest flight is $\val{errStep}$. The panel-size estimate
    is $\val{errMesh}$. The sampled clearance is $\val{refDmin}$, outside the
    shell $\val{contact}$. The image-force ratio reaches $\val{refImage}$.
  \item\label{st:rad} \emph{The radiation.} Simpson's rule on steps of $h/2$
    gives $\val{LowRad}$, $\val{AxisRad}$ and $\val{HighRad}$. The largest
    fraction of $T_0$ is $\val{radFracMax}$. Halving that step again changes
    the largest of those integrals by at most the relative amount
    $\val{HighRadErr}$. These integrals are not rows of the comparison.
\end{procedure}

\begin{result}[The reference flights]
All three ions pass the gate and enter the detector through its left side,
inside the cone. The level is solved.
\end{result}

\section{Comparison with the game}\label{sec:comparison}

\comparisontable

The table holds the outcome, the time and the end position of each reference
ray, and the same three for the empty placement. The tolerance on a state is
$\val{tolFlight}$. Ten times the panel-size estimate $\val{errMesh}$ lies
above $10^{-9}$, and it is larger than the step-halving estimate
$\val{errStep}$, so the tolerance is that product. The tolerance on a state
was fixed from these estimates before the comparison.

The largest absolute difference is the upward reference time,
$\val{cmpHighTDiff}$. It lies above the panel-size estimate and under the
tolerance. The downward reference time differs by $\val{cmpLowTDiff}$, also
above the panel-size estimate and under the tolerance. The outcomes agree.

The three reference abscissae agree exactly. Each is the detector's left edge.
The empty downward abscissa and the empty on-axis abscissa agree exactly, and
each is the world's right edge. The empty upward abscissa agrees exactly: the
ray meets the deflector's contact shell.

\section{Other solutions}\label{sec:others}

The empty placement in \code{placements.json} turns the supplies off and adds
no charge. The fixed charge and the grounded electrodes remain. It does not
solve the level.

The downward ray misses the gate and leaves through the right edge of the
world at $(\val{OffLowX}, \val{OffLowY})$, at $t = \val{OffLowT}$, at an angle
$\val{OffLowAng}$ degrees from $+x$. The on-axis ray does pass the gate, at
$(\val{OffAxisGX}, \val{OffAxisGY})$, and then leaves through the right edge
at $(\val{OffAxisX}, \val{OffAxisY})$, at $t = \val{OffAxisT}$, at
$\val{OffAxisAng}$ degrees. That ordinate is above the launch line. The
upward ray misses the gate and meets the deflector at
$(\val{OffHighX}, \val{OffHighY})$, at $t = \val{OffHighT}$, so the outcome is
\val{OffHighOutcome}. Its image-force ratio reaches $\val{OffHighImage}$.

The reference supplies and the added charge are what bring the three rays from
those three endings onto the detector, inside the cone. This document
integrates the reference and the empty placement.

\section{What the game computes}\label{sec:engine}

The electrodes are metal boxes, and a supply sets the potential of the
electrode it sits on (\physref{2.7}). The surface charge is the engine's
triangle collocation, the same boundary condition as
\xeqref{deflection_plates}{eq:collocation}, with the charges on the
right-hand side. This document uses the equal rectangles of the deflection
plates. The end states move by $\val{errMesh}$ when that edge doubles, and
the upward reference time in \S\ref{sec:comparison} lies above that estimate
and under ten times it.

The gate is \xeqref{two_stages}{eq:gate} and the cone is
\xeqref{injection}{eq:cone}. The ion's image force is omitted on both sides.
The bound sampled on the reference flights reaches $\val{refImage}$, under
the engine's limit $10^{-10}$. The flight in the engine is the adaptive
method of \physref{5} and \physref{6}, verified as in \physref{7}. Our step
estimate is $\val{errStep}$. Radiation is integrated and not applied to the
momentum (\xsecref{foundations}{sec:radiation}, \physref{8}).

\section{Further reading}\label{sec:reading}

The rectangle integral, the collocation and the flight energy are derived for
the deflection plates (\xeqref{deflection_plates}{eq:rectangle},
\xeqref{deflection_plates}{eq:collocation},
\xeqref{deflection_plates}{eq:flight-energy}). The aperture lens is derived
for Tune the lens (\xeqref{tune_the_lens}{eq:axis},
\xeqref{tune_the_lens}{eq:slow}). The gate is derived for Two stages
(\xeqref{two_stages}{eq:gate}). The cone is derived for Injection
(\xeqref{injection}{eq:cone}). The motion and the Runge--Kutta step are in
the Foundations (\xsecref{foundations}{sec:numerics}). The engine's
electrodes are \physref{2.7}.

\levelfileappendix

\end{document}
